\documentclass{article}
\usepackage{mathtools} % loads amsmath
\usepackage{amssymb}
\usepackage{amsthm}
\usepackage{hyperref}
\usepackage{cleveref} % must come after hyperref
\usepackage{thmtools} % \declaretheorem: lets \cref name shared-counter environments correctly

\declaretheorem[name=Theorem]{theorem}
\declaretheorem[name=Lemma, sibling=theorem]{lemma}
\declaretheorem[name=Definition, style=definition, sibling=theorem]{definition}

\title{Project Euler Problem 1: Multiples of 3 or 5}
\author{Brinhasavlin}
\date{\today}

\begin{document}
\maketitle

\section*{Problem}

If we list all the natural numbers below $10$ that are multiples of $3$ or $5$, we get $3, 5, 6$ and $9$. The sum of these multiples is $23$.

Find the sum of all the multiples of $3$ or $5$ below $1000$.

\section*{Proof}
Rather than testing every number below $n$, we sum the multiples of $3$ and of $5$ directly and correct for the numbers
counted twice. Throughout, we use the closed form for the sum of the first $k$ integers, proven by induction in
Problem 6:
\[
  T(k) = \sum_{j=1}^{k} j = \frac{k(k+1)}{2}.
\]

\begin{definition}\label{def:multsum}
  For $d, n \in \mathbb{N}$, let $\operatorname{S}_d(n)$ be the sum of the positive multiples of $d$ below $n$.
\end{definition}

\begin{lemma}\label{lem:multsum}
  For all $d, n \in \mathbb{N}$,
  \[
    \operatorname{S}_d(n) = d \cdot T\!\left(\left\lfloor \frac{n-1}{d} \right\rfloor\right).
  \]
\end{lemma}
\begin{proof}
  The positive multiples of $d$ are $jd$ for $j \ge 1$. Such a multiple is below $n$ exactly when $jd \le n - 1$, that
  is, when $j \le \frac{n-1}{d}$. Since $j$ is an integer, this is equivalent to $j \le k$, where
  $k = \left\lfloor \frac{n-1}{d} \right\rfloor$. Hence the multiples of $d$ below $n$ are $d, 2d, \dots, kd$, and
  \[
    \operatorname{S}_d(n) = \sum_{j=1}^{k} jd = d \sum_{j=1}^{k} j = d \cdot T(k).
  \]
  If $k = 0$, there are no such multiples, and both sides are $0$.
\end{proof}

\begin{definition}\label{def:lcm}
  For a finite nonempty set $D$ of positive integers, the \emph{least common multiple} $\operatorname{lcm}(D)$ is the
  smallest positive integer that is a multiple of every element of $D$.
\end{definition}

For this definition to make sense, a smallest such integer must exist. Let $P = \prod_{e \in D} e$ be the product of
all elements of $D$. For every $d \in D$,
\[
  P = d \cdot \prod_{e \in D \setminus \{d\}} e,
\]
where the second factor is a product of positive integers and is therefore a positive integer (it equals $1$ when
$D = \{d\}$, since an empty product is $1$). Hence $d \mid P$ for every $d \in D$, and $P > 0$ as a product of
positive integers. So the set of positive common multiples of $D$,
\[
  M = \{ m \in \mathbb{Z}_{>0} : d \mid m \text{ for every } d \in D \},
\]
contains $P$ and is nonempty. By the well-ordering principle, $M$ has a least element, and that element is
$\operatorname{lcm}(D)$. Hence $\operatorname{lcm}(D)$ is well defined.

\begin{lemma}\label{lem:multiple}
  A positive integer $x$ is a multiple of all elements of a finite nonempty set $D$ of positive integers if and only if
  it is also a multiple of $\operatorname{lcm}(D)$.
\end{lemma}
\begin{proof}
  Let $m = \operatorname{lcm}(D)$. We first show that if $x$ is a multiple of $m$, then it is a multiple of every
  element of $D$. Say $x = km$ for some integer $k$. Let $d \in D$. By \cref{def:lcm}, $d \mid m$, so there is an
  integer $p$ with $m = pd$. Then
  \[
    x = km = k(pd) = (kp)d,
  \]
  and $kp$ is an integer, so $d \mid x$. Since $d$ was arbitrary, $x$ is a multiple of every element of $D$.
  
  Conversely, we show that if $x$ is a multiple of every element of $D$, then it is necessarily a multiple of $m$ as
  well. Using the set $M$ defined right after \Cref{def:lcm}, we know that $m$ is the smallest multiple of every
  element of $D$ and that $m$ is greater than $0$. We start by dividing $x$ by $m$:
  \begin{align*}
    x &= qm + r, \qquad 0 \le r < m, \\
    r &= x - qm.
  \end{align*}
  For an arbitrary $d$ in $D$, we know that $d \mid x$ by the hypothesis. Since $m \in M$, we also know that
  $d \mid m$, implying that $d \mid qm$. So there are integers $a$ and $b$ with $x = ad$ and $qm = bd$, and
  \[
    r = x - qm = ad - bd = (a - b)d.
  \]
  Therefore, $d \mid r$, and since $d$ was arbitrary, every element of $D$ divides $r$. If we suppose that $r > 0$, then $r$ would be a member of $M$ that is smaller than $m$. But $m$ is
  defined to be the smallest member of $M$, so $r = 0$. This means that $x = qm$, where $q$ is an integer, making $x$
  a multiple of $m$.
\end{proof}

With only two divisors, it is easy to see that subtracting the multiples of $15$ once fixes the double counting. The
same idea works for any set of divisors, but with more of them the overlaps are harder to track by hand, so we prove
the general statement by looking at how many times each number is counted.

\begin{lemma}\label{lem:incexc}
  Let $D$ be a finite nonempty set of positive integers and let $n \in \mathbb{N}$. The sum of the positive integers
  below $n$ that are multiples of at least one element of $D$ is
  \[
    \sum_{\emptyset \ne A \subseteq D} {(-1)}^{|A|+1} \operatorname{S}_{\operatorname{lcm}(A)}(n).
  \]
\end{lemma}
\begin{proof}
  The sum we want adds each $x$ below $n$ that is a multiple of some element of $D$ exactly once, and leaves out every
  other $x$. We rewrite the alternating sum in the statement as a single sum over $x$, where each $x$ is multiplied by a
  coefficient that we call its \emph{weight}. We then show that the weight of $x$ is $1$ if $x$ is a multiple of some
  element of $D$ and $0$ otherwise, so the two sums agree.

  For each integer $x$ with $1 \le x < n$, let $D_x = \{d \in D : d \mid x\}$ be the set of divisors in $D$ that
  divide $x$, and let $m_x = |D_x|$.

  Let $A \subseteq D$ be nonempty. Then $A$ is a finite nonempty set of positive integers, so by \cref{lem:multiple},
  a positive integer is a multiple of every element of $A$ if and only if it is a multiple of
  $\operatorname{lcm}(A)$. So $\operatorname{S}_{\operatorname{lcm}(A)}(n)$ is the sum of those $x < n$
  that are divisible by every element of $A$, that is, those with $A \subseteq D_x$. Writing $[P]$ for $1$ if the
  statement $P$ is true and $0$ otherwise, we can sum over every candidate below $n$ and keep only those that qualify:
  \[
    \operatorname{S}_{\operatorname{lcm}(A)}(n) = \sum_{x=1}^{n-1} x \, [A \subseteq D_x].
  \]
  Substituting this, swapping the order of the two finite sums, and factoring out $x$, which does not depend on $A$,
  \begin{align*}
    \sum_{\emptyset \ne A \subseteq D} {(-1)}^{|A|+1} \operatorname{S}_{\operatorname{lcm}(A)}(n)
      &= \sum_{\emptyset \ne A \subseteq D} {(-1)}^{|A|+1} \sum_{x=1}^{n-1} x \, [A \subseteq D_x] \\
      &= \sum_{x=1}^{n-1} x \sum_{\emptyset \ne A \subseteq D} {(-1)}^{|A|+1} [A \subseteq D_x] \\
      &= \sum_{x=1}^{n-1} x \sum_{\emptyset \ne A \subseteq D_x} {(-1)}^{|A|+1}.
  \end{align*}
  The last step uses the bracket to drop the subsets that are not contained in $D_x$.
  The inner sum is the weight of $x$. To compute it, we group the nonempty subsets of $D_x$ by size. For each $j$ with
  $1 \le j \le m_x$, there are $\binom{m_x}{j}$ subsets of size $j$, and each contributes ${(-1)}^{j+1}$. Hence
  \[
    \sum_{\emptyset \ne A \subseteq D_x} {(-1)}^{|A|+1} = \sum_{j=1}^{m_x} {(-1)}^{j+1} \binom{m_x}{j} = w(m_x),
  \]
  where for any integer $k \ge 0$ we define
  \[
    w(k) = \sum_{j=1}^{k} {(-1)}^{j+1} \binom{k}{j}.
  \]
  If $k = 0$, the sum is empty and $w(0) = 0$. If $k \ge 1$, then
  \begin{align*}
    0 = {(1 - 1)}^k
      &= \sum_{j=0}^{k} \binom{k}{j} 1^{k-j} {(-1)}^{j} && \text{by the binomial theorem} \\
      &= \sum_{j=0}^{k} \binom{k}{j} {(-1)}^{j} && \text{since $1^{k-j} = 1$} \\
      &= 1 + \sum_{j=1}^{k} \binom{k}{j} {(-1)}^{j} && \text{pulling out the $j = 0$ term} \\
      &= 1 - \sum_{j=1}^{k} \binom{k}{j} {(-1)}^{j+1} && \text{since ${(-1)}^{j} = -{(-1)}^{j+1}$} \\
      &= 1 - w(k),
  \end{align*}
  so $w(k) = 1$. Thus $w(m_x) = 1$ when $m_x \ge 1$, that is, when $x$ is a multiple of some element of $D$, and
  $w(m_x) = 0$ otherwise. Therefore
  \[
    \sum_{\emptyset \ne A \subseteq D} {(-1)}^{|A|+1} \operatorname{S}_{\operatorname{lcm}(A)}(n)
      = \sum_{x=1}^{n-1} x \, w(m_x)
      = \sum_{\substack{1 \le x < n \\ m_x \ge 1}} x,
  \]
  which is the sum of the positive integers below $n$ that are multiples of at least one element of $D$.
\end{proof}

\begin{theorem}\label{thm:sum}
  For all $n \in \mathbb{N}$, the sum of the positive integers below $n$ that are multiples of $3$ or $5$ is
  \[
    \operatorname{S}_3(n) + \operatorname{S}_5(n) - \operatorname{S}_{15}(n).
  \]
\end{theorem}
\begin{proof}
  Apply \cref{lem:incexc} with $D = \{3, 5\}$. The nonempty subsets are $\{3\}$ and $\{5\}$, each of size $1$, and
  $\{3, 5\}$, of size $2$. A single-element set $\{d\}$ has $\operatorname{lcm}(\{d\}) = d$, since $d$ is the smallest
  positive multiple of $d$. For $\{3, 5\}$, the number $15$ is a positive multiple of both $3$ and $5$. Any smaller
  common multiple would be a positive multiple of $5$ below $15$, so it would be $5$ or $10$, and $3$ divides neither.
  Hence $\operatorname{lcm}(\{3, 5\}) = 15$, and the three terms are $\operatorname{S}_3(n)$ and
  $\operatorname{S}_5(n)$ with sign $+$, and $\operatorname{S}_{15}(n)$ with sign $-$.
\end{proof}

\section*{Answer}
As a check, the example from the problem statement ($n = 10$) gives
\[
  3 \cdot T(3) + 5 \cdot T(1) - 15 \cdot T(0) = 18 + 5 - 0 = 23.
\]
For $n = 1000$, we have $\lfloor 999/3 \rfloor = 333$, $\lfloor 999/5 \rfloor = 199$, and $\lfloor 999/15 \rfloor = 66$.
By \cref{thm:sum,lem:multsum},
\begin{align*}
  \operatorname{S}_3(1000) + \operatorname{S}_5(1000) - \operatorname{S}_{15}(1000)
    &= 3 \cdot T(333) + 5 \cdot T(199) - 15 \cdot T(66) \\
    &= 3 \cdot 55611 + 5 \cdot 19900 - 15 \cdot 2211 \\
    &= 166833 + 99500 - 33165 \\
    &= 233168.
\end{align*}

\end{document}
